\documentclass[11pt,letterpaper]{article}

% ---------------------------------------------------------------
% FoodBridge -- Milestone 3 (CIS 5500, Group 25)
% Monochrome layout: black, white and shades of gray only.
% Compiles with pdfLaTeX on Overleaf.
% ---------------------------------------------------------------
\usepackage[T1]{fontenc}
\usepackage[utf8]{inputenc}
\usepackage{lmodern}
\usepackage[margin=1in,headheight=14pt]{geometry}
\usepackage{xcolor}
\usepackage{amssymb}
\usepackage{array}
\usepackage{booktabs}
\usepackage{tabularx}
\usepackage{titlesec}
\usepackage{fancyhdr}
\usepackage{listings}
\usepackage{upquote}
\usepackage{microtype}
\usepackage[hidelinks]{hyperref}

% ---------- grays ----------
\definecolor{ruleGray}{gray}{0.55}
\definecolor{codeBg}{gray}{0.96}
\definecolor{codeFrame}{gray}{0.70}
\definecolor{commentGray}{gray}{0.40}
\definecolor{stringGray}{gray}{0.25}
\definecolor{metaGray}{gray}{0.35}

% ---------- paragraphs ----------
\setlength{\parindent}{0pt}
\setlength{\parskip}{6pt}

% ---------- headings ----------
\titleformat{\section}{\Large\bfseries}{}{0pt}{}[{\color{ruleGray}\titlerule[0.6pt]}]
\titlespacing*{\section}{0pt}{18pt}{8pt}
\titleformat{\subsection}{\large\bfseries}{}{0pt}{}
\titlespacing*{\subsection}{0pt}{14pt}{4pt}

% ---------- header / footer ----------
\pagestyle{fancy}
\fancyhf{}
\fancyhead[L]{\small\color{metaGray}FoodBridge --- Milestone 3}
\fancyhead[R]{\small\color{metaGray}CIS 5500 $\cdot$ Group 25}
\fancyfoot[C]{\small\color{metaGray}\thepage}
\renewcommand{\headrulewidth}{0.4pt}
\renewcommand{\headrule}{{\color{ruleGray}\hrule width\headwidth height\headrulewidth}}

% ---------- SQL listings ----------
\lstdefinestyle{sqlmono}{
  language=SQL,
  basicstyle=\ttfamily\footnotesize\color{black},
  keywordstyle=\bfseries\color{black},
  commentstyle=\itshape\color{commentGray},
  stringstyle=\color{stringGray},
  backgroundcolor=\color{codeBg},
  frame=single,
  rulecolor=\color{codeFrame},
  framesep=5pt,
  xleftmargin=5pt,
  xrightmargin=5pt,
  aboveskip=6pt,
  belowskip=10pt,
  columns=fullflexible,
  keepspaces=true,
  showstringspaces=false,
  breaklines=true,
  breakatwhitespace=true,
  postbreak=\mbox{\textcolor{commentGray}{$\hookrightarrow$}\space},
  upquote=true,
  tabsize=4,
  morekeywords={WITH, FILTER, ILIKE, RANK, OVER, INTERVAL, TIMESTAMP, COALESCE, NULLIF,
                GREATEST, ROUND, LIMIT, FULL, OUTER, EXISTS, date_trunc, to_char}
}
\lstset{style=sqlmono}

\renewcommand{\arraystretch}{1.25}
\newcolumntype{C}{>{\centering\arraybackslash}X}

\begin{document}

% ================= Title =================
\begin{center}
  {\LARGE\bfseries FoodBridge -- Milestone 3:\\[4pt] Database Population \& SQL Queries\par}
  \vspace{8pt}
  {\normalsize\color{metaGray} November 10, 2026 \quad$\cdot$\quad Sumit Pradeep Kolpe\par}
  \vspace{4pt}
  {\color{ruleGray}\rule{\linewidth}{0.8pt}}
\end{center}

The FoodBridge database is fully populated on AWS RDS (PostgreSQL): 5 tables, 469,440 rows. Below are guest credentials and 12 SQL queries (6 simple, 6 complex) that run on it.

% ================= Team =================
\section*{Team}
CIS 5500 -- Database and Information Systems, Fall 2026 $\cdot$ Group 25

\begin{center}
\begin{tabularx}{\linewidth}{>{\raggedright\arraybackslash}X l l}
\toprule
\textbf{Name} & \textbf{Email} & \textbf{GitHub} \\
\midrule
Sumit Pradeep Kolpe & \texttt{sumit1@engineering.upenn.edu} & \texttt{sumit-kolpe} \\
Urav Dalal & \texttt{urav@engineering.upenn.edu} & \texttt{Urav-Dalal} \\
Jigar Shah & \texttt{jigars@engineering.upenn.edu} & \texttt{shahjigar556} \\
Rushabh Maru & \texttt{rpmaru@engineering.upenn.edu} & \texttt{RushabhM03} \\
\bottomrule
\end{tabularx}
\end{center}

% ================= Access =================
\section*{Database access}
Use the read-only guest account below. It can run \texttt{SELECT} on every table but cannot modify data, and each statement is limited to 10 minutes.

\begin{center}
\begin{tabularx}{0.75\linewidth}{>{\bfseries}l X}
\toprule
Setting & \textbf{Value} \\
\midrule
Engine & PostgreSQL on AWS RDS \\
Host (endpoint) & \texttt{<RDS endpoint>} \\
Port & \texttt{5432} \\
Database & \texttt{foodbridge} \\
Username & \texttt{fb\_guest} \\
Password & \texttt{<guest password>} \\
\bottomrule
\end{tabularx}
\end{center}

Connect with psql:
\begin{lstlisting}[language=bash]
psql -h <endpoint> -p 5432 -U fb_guest -d foodbridge
\end{lstlisting}

Or in DataGrip / pgAdmin: new PostgreSQL data source with the host, port, database, user and password above.

The instance runs in AWS Academy Learner Lab. If it is unreachable, the lab session may be stopped; please email any team member and we will start it.

% ================= Contents =================
\section*{Database contents}
The schema is the Milestone 2 DDL (3NF), loaded with cleaned real and synthetic data. All primary-key, foreign-key and \texttt{CHECK} constraints are enforced, and 41 validation checks passed before loading.

\begin{center}
\begin{tabularx}{\linewidth}{l r X}
\toprule
\textbf{Table} & \textbf{Rows} & \textbf{Source} \\
\midrule
Organization & 1,000 & 346 real Philadelphia food sites (ids 1--346, after cleaning and entity resolution of 404 raw records) + 654 synthetic donors and recipients \\
FoodCategory & 25 & Curated reference list \\
FoodDonation & 150,000 & Synthetic, Jan 2025 -- Nov 2026 \\
RecipientRequest & 120,000 & Synthetic, Jan 2025 -- Nov 2026 \\
Allocation & 198,415 & Simulated matching of donations to requests \\
\bottomrule
\end{tabularx}
\end{center}

The data is a snapshot taken at 2026-11-10 12:00:00. Queries use that timestamp as ``now'' (an \texttt{as\_of} parameter in the API) instead of \texttt{NOW()}. Literal ids, ZIP codes and dates in the queries are example parameter values.

The full SQL is also in our repository at \texttt{sql/queries.sql}.

% ================= Simple queries =================
\section*{Simple queries (Q1--Q6)}
Each simple query backs one feature of a page in the application.

\subsection*{Q1. Organization directory search --- Organizations page}
Lists active recipient organizations in a given ZIP code whose name matches a search term, in alphabetical order.

\begin{lstlisting}
SELECT o.organization_id, o.organization_name, o.organization_type,
       o.address, o.zip_code, o.contact_information
FROM Organization o
WHERE o.status = 'active'
  AND o.organization_type IN ('recipient', 'both')
  AND o.zip_code = '19140'
  AND o.organization_name ILIKE '%church%'
ORDER BY o.organization_name;
\end{lstlisting}

\subsection*{Q2. Claimable donations feed --- Food Donations page}
Shows donations that can still be picked up at the reference time, with donor, category, storage type and servings not yet allocated, soonest pickup deadline first.

\begin{lstlisting}
SELECT d.donation_id, o.organization_name AS donor, c.category_name, c.storage_type,
       d.quantity,
       d.quantity - COALESCE(SUM(a.quantity_allocated)
                             FILTER (WHERE a.status <> 'cancelled'), 0) AS remaining_servings,
       d.pickup_deadline, d.expires_at
FROM FoodDonation d
JOIN Organization o ON o.organization_id = d.donor_org_id
JOIN FoodCategory c ON c.category_id = d.category_id
LEFT JOIN Allocation a ON a.donation_id = d.donation_id
WHERE d.status IN ('available', 'partially_allocated')
  AND d.pickup_deadline > TIMESTAMP '2026-11-10 12:00:00'
GROUP BY d.donation_id, o.organization_name, c.category_name, c.storage_type
ORDER BY d.pickup_deadline
LIMIT 50;
\end{lstlisting}

\subsection*{Q3. Donor profile summary --- Organization detail page}
For one donor organization, returns the number of donations, total servings and average donation size in each food category.

\begin{lstlisting}
SELECT c.category_name,
       COUNT(*)                  AS donations,
       SUM(d.quantity)           AS servings_donated,
       ROUND(AVG(d.quantity), 1) AS avg_servings
FROM FoodDonation d
JOIN FoodCategory c ON c.category_id = d.category_id
WHERE d.donor_org_id = 400
GROUP BY c.category_name
ORDER BY servings_donated DESC;
\end{lstlisting}

\subsection*{Q4. Request tracking --- Recipient Requests page}
Lists every allocation made against one request: the supplying donor, servings, and pickup/delivery status and times.

\begin{lstlisting}
SELECT a.allocation_id, o.organization_name AS donor, a.quantity_allocated,
       a.status, a.allocated_at, a.picked_up_at, a.delivered_at
FROM Allocation a
JOIN FoodDonation d ON d.donation_id = a.donation_id
JOIN Organization o ON o.organization_id = d.donor_org_id
WHERE a.request_id = 100006
ORDER BY a.allocated_at;
\end{lstlisting}

\subsection*{Q5. Platform KPIs --- Dashboard}
Returns the dashboard headline numbers in one row: allocations, delivered allocations and servings, servings in transit, and cancellations.

\begin{lstlisting}
SELECT COUNT(*)                                                       AS total_allocations,
       COUNT(*) FILTER (WHERE status = 'delivered')                   AS delivered_allocations,
       SUM(quantity_allocated) FILTER (WHERE status = 'delivered')    AS servings_delivered,
       SUM(quantity_allocated) FILTER (WHERE status IN ('pending', 'confirmed', 'picked_up'))
                                                                      AS servings_in_transit,
       COUNT(*) FILTER (WHERE status = 'cancelled')                   AS cancelled_allocations
FROM Allocation;
\end{lstlisting}

\subsection*{Q6. Category demand in the last 30 days --- Dashboard / Requests page}
Counts requests and requested servings per food category over the 30 days before the reference time, busiest categories first.

\begin{lstlisting}
SELECT c.category_name, c.storage_type,
       COUNT(r.request_id)       AS requests,
       SUM(r.requested_quantity) AS servings_requested
FROM FoodCategory c
JOIN RecipientRequest r ON r.category_id = c.category_id
WHERE r.requested_at >= TIMESTAMP '2026-11-10 12:00:00' - INTERVAL '30 days'
  AND r.status <> 'cancelled'
GROUP BY c.category_name, c.storage_type
ORDER BY servings_requested DESC;
\end{lstlisting}


% ================= Complex queries =================
\section*{Complex queries (Q7--Q12)}
Each complex query combines data from several tables, including the real Philadelphia sites and the synthetic operational data. Q10 and Q12 are the long-running queries we will optimize; Q11 is a universal (``for all'') check.

\begin{center}
{\small
\begin{tabularx}{\linewidth}{>{\raggedright\arraybackslash}p{0.30\linewidth} C C C >{\centering\arraybackslash}p{0.18\linewidth}}
\toprule
\textbf{Query} & \textbf{Multiple joins} & \textbf{Subqueries / CTEs} & \textbf{Aggregation} & \textbf{Exists / for-all check} \\
\midrule
Q7 Monthly supply, demand and fulfilment & \checkmark & \checkmark & \checkmark & \\
Q8 Donor performance benchmarking & \checkmark & \checkmark & \checkmark & \\
Q9 Geographic distribution gaps & \checkmark & \checkmark & \checkmark & \\
Q10 Avoidable food loss & \checkmark & \checkmark & \checkmark & \texttt{EXISTS} \\
Q11 Fully served recipients & \checkmark & \checkmark & \checkmark & \texttt{NOT EXISTS} /\newline \texttt{NOT EXISTS}\newline (for all) \\
Q12 Structurally unmet demand & \checkmark & \checkmark & \checkmark & \\
\bottomrule
\end{tabularx}
}
\end{center}

\subsection*{Q7. Monthly supply, demand and fulfilment by category --- Analytics page}
For each month and food category, compares servings donated, servings requested and servings actually delivered, and the share of demand that was met.

\begin{lstlisting}
WITH supply AS (
    SELECT date_trunc('month', donated_at) AS month, category_id,
           SUM(quantity) AS donated
    FROM FoodDonation
    WHERE status <> 'cancelled'
    GROUP BY 1, 2
), demand AS (
    SELECT date_trunc('month', requested_at) AS month, category_id,
           SUM(requested_quantity) AS requested
    FROM RecipientRequest
    WHERE status <> 'cancelled'
    GROUP BY 1, 2
), delivered AS (
    SELECT date_trunc('month', r.requested_at) AS month, r.category_id,
           SUM(a.quantity_allocated) AS delivered
    FROM Allocation a
    JOIN RecipientRequest r ON r.request_id = a.request_id
    WHERE a.status = 'delivered'
    GROUP BY 1, 2
)
SELECT to_char(COALESCE(s.month, dm.month), 'YYYY-MM')       AS month,
       c.category_name,
       COALESCE(s.donated, 0)                                 AS servings_donated,
       COALESCE(dm.requested, 0)                              AS servings_requested,
       COALESCE(dl.delivered, 0)                              AS servings_delivered,
       ROUND(100.0 * COALESCE(dl.delivered, 0) / NULLIF(dm.requested, 0), 1) AS demand_met_pct
FROM supply s
FULL OUTER JOIN demand dm ON dm.month = s.month AND dm.category_id = s.category_id
LEFT JOIN delivered dl
       ON dl.month = COALESCE(s.month, dm.month)
      AND dl.category_id = COALESCE(s.category_id, dm.category_id)
JOIN FoodCategory c ON c.category_id = COALESCE(s.category_id, dm.category_id)
ORDER BY month, c.category_name;
\end{lstlisting}

\subsection*{Q8. Donor performance benchmarking --- Analytics page (donor leaderboard)}
For donors with at least 20 donations, ranks them by servings that reached a recipient, and compares each donor's delivery rate with the platform-wide rate.

\begin{lstlisting}
WITH donation_outcome AS (
    SELECT d.donation_id, d.donor_org_id, d.quantity,
           COALESCE(SUM(a.quantity_allocated) FILTER (WHERE a.status = 'delivered'), 0) AS delivered
    FROM FoodDonation d
    LEFT JOIN Allocation a ON a.donation_id = d.donation_id
    WHERE d.status <> 'cancelled'
    GROUP BY d.donation_id, d.donor_org_id, d.quantity
), donor_stats AS (
    SELECT o.organization_id, o.organization_name, o.zip_code,
           COUNT(*)            AS donations,
           SUM(x.quantity)     AS servings_donated,
           SUM(x.delivered)    AS servings_delivered,
           ROUND(100.0 * SUM(x.delivered) / SUM(x.quantity), 1) AS delivery_rate_pct
    FROM donation_outcome x
    JOIN Organization o ON o.organization_id = x.donor_org_id
    GROUP BY o.organization_id, o.organization_name, o.zip_code
    HAVING COUNT(*) >= 20
)
SELECT RANK() OVER (ORDER BY servings_delivered DESC) AS rank_by_delivered,
       organization_name, zip_code, donations, servings_donated, servings_delivered,
       delivery_rate_pct,
       ROUND(delivery_rate_pct - (SELECT 100.0 * SUM(delivered) / SUM(quantity)
                                  FROM donation_outcome), 1) AS vs_platform_pct_points
FROM donor_stats
ORDER BY rank_by_delivered
LIMIT 25;
\end{lstlisting}

\subsection*{Q9. Geographic distribution gaps --- Analytics page (map)}
For each ZIP code, computes undelivered servings of requests, the number of active recipient sites and the servings donated by local donors, then ranks ZIP codes by unmet need per site.

\begin{lstlisting}
WITH request_outstanding AS (
    SELECT r.request_id, r.recipient_org_id,
           r.requested_quantity
             - COALESCE(SUM(a.quantity_allocated) FILTER (WHERE a.status = 'delivered'), 0) AS outstanding
    FROM RecipientRequest r
    LEFT JOIN Allocation a ON a.request_id = r.request_id
    WHERE r.status IN ('open', 'partially_fulfilled', 'expired')
    GROUP BY r.request_id, r.recipient_org_id, r.requested_quantity
), zip_need AS (
    SELECT o.zip_code, SUM(ro.outstanding) AS unmet_servings,
           COUNT(*) FILTER (WHERE ro.outstanding > 0) AS unmet_requests
    FROM request_outstanding ro
    JOIN Organization o ON o.organization_id = ro.recipient_org_id
    WHERE o.zip_code IS NOT NULL
    GROUP BY o.zip_code
), zip_sites AS (
    SELECT zip_code, COUNT(*) AS active_recipient_sites
    FROM Organization
    WHERE status = 'active' AND organization_type IN ('recipient', 'both')
    GROUP BY zip_code
), zip_supply AS (
    SELECT o.zip_code, SUM(d.quantity) AS servings_donated_locally
    FROM FoodDonation d
    JOIN Organization o ON o.organization_id = d.donor_org_id
    GROUP BY o.zip_code
)
SELECT n.zip_code, n.unmet_servings, n.unmet_requests,
       COALESCE(s.active_recipient_sites, 0)          AS active_recipient_sites,
       COALESCE(sp.servings_donated_locally, 0)        AS servings_donated_locally,
       ROUND(n.unmet_servings::numeric / GREATEST(COALESCE(s.active_recipient_sites, 0), 1), 0)
                                                       AS unmet_per_site,
       RANK() OVER (ORDER BY n.unmet_servings::numeric
                    / GREATEST(COALESCE(s.active_recipient_sites, 0), 1) DESC) AS gap_rank
FROM zip_need n
LEFT JOIN zip_sites s  ON s.zip_code = n.zip_code
LEFT JOIN zip_supply sp ON sp.zip_code = n.zip_code
ORDER BY gap_rank;
\end{lstlisting}

\subsection*{Q10. Avoidable food loss --- Analytics page (wastage)}
Finds donations from the last 14 days that expired with servings left over even though a compatible request (same category, open during the pickup window, still short of servings) existed, and totals that avoidable loss per category and donor.

\begin{lstlisting}
WITH expired_leftover AS (
    SELECT d.donation_id, d.donor_org_id, d.category_id, d.donated_at, d.pickup_deadline,
           d.quantity - COALESCE((SELECT SUM(a.quantity_allocated)
                                  FROM Allocation a
                                  WHERE a.donation_id = d.donation_id
                                    AND a.status <> 'cancelled'), 0) AS leftover
    FROM FoodDonation d
    WHERE d.status = 'expired'
      AND d.donated_at >= TIMESTAMP '2026-11-10 12:00:00' - INTERVAL '14 days'
)
SELECT c.category_name, o.organization_name AS donor,
       COUNT(*)        AS avoidable_expired_donations,
       SUM(e.leftover) AS avoidable_lost_servings
FROM expired_leftover e
JOIN FoodCategory c ON c.category_id = e.category_id
JOIN Organization o ON o.organization_id = e.donor_org_id
WHERE e.leftover > 0
  AND EXISTS (
        SELECT 1
        FROM RecipientRequest r
        WHERE r.category_id = e.category_id
          AND r.status <> 'cancelled'
          AND r.requested_at >= TIMESTAMP '2026-11-10 12:00:00' - INTERVAL '14 days'
          AND r.requested_at <= e.pickup_deadline
          AND r.needed_by    >= e.donated_at
          AND r.requested_quantity > COALESCE((SELECT SUM(a2.quantity_allocated)
                                               FROM Allocation a2
                                               WHERE a2.request_id = r.request_id
                                                 AND a2.status <> 'cancelled'), 0))
GROUP BY c.category_name, o.organization_name
ORDER BY avoidable_lost_servings DESC
LIMIT 25;
\end{lstlisting}

\subsection*{Q11. Fully served recipients --- Analytics page (recipient coverage)}
Finds active recipient organizations that, in the last 180 days, received at least one delivery in every food category they requested, with their categories, requests and servings received. This is relational division written as a double \texttt{NOT EXISTS}.

\begin{lstlisting}
SELECT o.organization_id, o.organization_name, o.zip_code,
       COUNT(DISTINCT r.category_id) AS categories_requested,
       COUNT(DISTINCT r.request_id)  AS requests,
       (SELECT SUM(a.quantity_allocated)
        FROM Allocation a
        JOIN RecipientRequest r3 ON r3.request_id = a.request_id
        WHERE r3.recipient_org_id = o.organization_id
          AND a.status = 'delivered'
          AND r3.requested_at >= TIMESTAMP '2026-11-10 12:00:00' - INTERVAL '180 days') AS servings_received
FROM Organization o
JOIN RecipientRequest r
  ON r.recipient_org_id = o.organization_id
 AND r.status <> 'cancelled'
 AND r.requested_at >= TIMESTAMP '2026-11-10 12:00:00' - INTERVAL '180 days'
WHERE o.status = 'active'
  AND o.organization_type IN ('recipient', 'both')
  AND NOT EXISTS (                       -- there is no requested category ...
        SELECT 1
        FROM RecipientRequest r1
        WHERE r1.recipient_org_id = o.organization_id
          AND r1.status <> 'cancelled'
          AND r1.requested_at >= TIMESTAMP '2026-11-10 12:00:00' - INTERVAL '180 days'
          AND NOT EXISTS (               -- ... that never received a delivery
                SELECT 1
                FROM Allocation a
                JOIN RecipientRequest r2 ON r2.request_id = a.request_id
                WHERE r2.recipient_org_id = o.organization_id
                  AND r2.category_id = r1.category_id
                  AND r2.requested_at >= TIMESTAMP '2026-11-10 12:00:00' - INTERVAL '180 days'
                  AND a.status = 'delivered'))
GROUP BY o.organization_id, o.organization_name, o.zip_code
ORDER BY categories_requested DESC, servings_received DESC;
\end{lstlisting}

\subsection*{Q12. Structurally unmet demand --- Recipient Requests page (``at risk'') and Analytics page}
Finds requests from the last 14 days whose outstanding servings exceed all remaining servings of compatible donations (same category, posted before the deadline, still collectable after the request), and reports per category how many such requests exist, how many recipients they affect and the total shortfall.

\begin{lstlisting}
WITH donation_remaining AS (
    SELECT d.donation_id, d.category_id, d.donated_at, d.pickup_deadline,
           d.quantity - COALESCE(SUM(a.quantity_allocated)
                                 FILTER (WHERE a.status <> 'cancelled'), 0) AS remaining
    FROM FoodDonation d
    LEFT JOIN Allocation a ON a.donation_id = d.donation_id
    WHERE d.status <> 'cancelled'
    GROUP BY d.donation_id, d.category_id, d.donated_at, d.pickup_deadline, d.quantity
), request_outstanding AS (
    SELECT r.request_id, r.recipient_org_id, r.category_id, r.requested_at, r.needed_by,
           r.requested_quantity - COALESCE(SUM(a.quantity_allocated)
                                           FILTER (WHERE a.status <> 'cancelled'), 0) AS outstanding
    FROM RecipientRequest r
    LEFT JOIN Allocation a ON a.request_id = r.request_id
    WHERE r.status IN ('open', 'partially_fulfilled', 'expired')
      AND r.requested_at >= TIMESTAMP '2026-11-10 12:00:00' - INTERVAL '14 days'
    GROUP BY r.request_id, r.recipient_org_id, r.category_id, r.requested_at, r.needed_by,
             r.requested_quantity
)
SELECT c.category_name,
       COUNT(*)                                    AS unmeetable_requests,
       COUNT(DISTINCT ro.recipient_org_id)          AS recipients_affected,
       SUM(ro.outstanding - ro.compatible_supply)   AS total_shortfall_servings
FROM (
    SELECT q.*,
           (SELECT COALESCE(SUM(dr.remaining), 0)
            FROM donation_remaining dr
            WHERE dr.category_id = q.category_id
              AND dr.donated_at <= q.needed_by
              AND dr.pickup_deadline >= q.requested_at) AS compatible_supply
    FROM request_outstanding q
    WHERE q.outstanding > 0
) ro
JOIN FoodCategory c ON c.category_id = ro.category_id
WHERE ro.outstanding > ro.compatible_supply
GROUP BY c.category_name
HAVING SUM(ro.outstanding - ro.compatible_supply) > 0
ORDER BY total_shortfall_servings DESC;
\end{lstlisting}


\end{document}
